\documentclass[11pt,a4paper]{article}
\usepackage[margin=24mm]{geometry}
\usepackage[T1]{fontenc}
\usepackage{lmodern,microtype,amsmath,amssymb,booktabs,longtable,array,enumitem,xcolor,graphicx}
\usepackage{tikz}
\usepackage{fancyhdr,needspace}
\let\originalsection\section
\renewcommand{\section}{\Needspace{7\baselineskip}\originalsection}
\definecolor{navy}{HTML}{17324D}
\definecolor{teal}{HTML}{087E8B}
\usepackage[colorlinks=true,linkcolor=navy,urlcolor=teal,citecolor=teal,bookmarksnumbered=true]{hyperref}
\hypersetup{pdftitle={Nuclear Medicine: Foundations, Activity and Decay},pdfauthor={Joaquin del Pilar}}
\setlength{\parindent}{0pt}
\setlength{\parskip}{6pt}
\setlength{\emergencystretch}{3em}
\setlist{itemsep=3pt,topsep=4pt,leftmargin=*}
\renewcommand{\arraystretch}{1.2}
\newcolumntype{P}[1]{>{\raggedright\arraybackslash}p{#1}}
\pagestyle{fancy}\fancyhf{}
\fancyhead[L]{\small\textcolor{navy}{NUCLEAR MEDICINE}}
\fancyhead[R]{\small Foundations and decay}
\fancyfoot[C]{\thepage}
\setlength{\headheight}{15pt}
\newcommand{\nucl}[3]{\ensuremath{{}^{#1}_{#2}\mathrm{#3}}}
\setcounter{tocdepth}{2}
\let\originalsubsection\subsection
\renewcommand{\subsection}{\Needspace{8\baselineskip}\originalsubsection}
\begin{document}
\hypersetup{pageanchor=false}
\begin{titlepage}
\vspace*{24mm}
{\color{teal}\large NUC MD 01\par}
\vspace{8mm}
{\color{navy}\Huge\bfseries Nuclear Medicine\par}
\vspace{6mm}
{\LARGE Foundations, Activity\par and Radioactive Decay\par}
\vspace{14mm}
{\large Lecture Notes\par}
\vfill
\end{titlepage}
\hypersetup{pageanchor=true}
{\small\tableofcontents}
\clearpage
\section{Physics Foundations}
Nuclear medicine is built on the production, utilization, measurement, and safe handling of \textbf{radionuclides}. The \textbf{medical physicist} is the bridge between theory, technology, and clinical application.

The course emphasizes a \textbf{physics-based foundation}, building on atomic and nuclear topics from Radiation Physics I. Assignments primarily serve \textbf{self-assessment} and may be discussed in subsequent meetings, including board work. The approach is \textbf{algebra-based}, with calculus appearing in guided examples.

\subsection{Radioactivity and Practical Work}
Radioactivity decays continuously. Transit times reduce the activity of radionuclides delivered to distant facilities and create logistical challenges.

Accurate activity measurements begin with a \textbf{dose calibrator}. Radiation safety emphasizes \textbf{ALARA (As Low As Reasonably Achievable)} and protection through \textbf{time, distance, and shielding}.

\subsection{Atomic and Nuclear Structure}

An \textbf{atom} is the smallest unit retaining chemical identity; atoms join to form molecules. Electrons surround a central \textbf{nucleus} composed of \textbf{protons} and \textbf{neutrons}. The number of protons defines the element.

\begin{itemize}
\item \textbf{Quantum numbers} describe electronic states:
\begin{itemize}
\item \textbf{Principal quantum number (n)} specifies shell/energy level.
\item \textbf{Orbital quantum number (l)} specifies orbital shape (s, p, d, f, g).
\item \textbf{Magnetic quantum number (m)} specifies spatial orientation in a magnetic field.
\item \textbf{Spin magnetic quantum number (m\ensuremath{{}_{s}})} specifies intrinsic electron spin (+1/2 or -1/2).
\end{itemize}
\end{itemize}
The three spatial orientations (px, py, pz) apply specifically to the \textbf{p subshell} (l = 1). Other subshells have different numbers of m-states (s: 1; d: 5; f: 7), consistent with 2l + 1 degeneracy.

Protons and neutrons are composite particles made of \textbf{quarks}:

\begin{itemize}
\item \textbf{Up quark}: charge +2/3
\item \textbf{Down quark}: charge -1/3
\item \textbf{Proton (uud)}: (+2/3) + (+2/3) + (-1/3) = \textbf{+1}
\item \textbf{Neutron (udd)}: (+2/3) + (-1/3) + (-1/3) = \textbf{0}
\end{itemize}
This quark composition explains the charges of nucleons.


\clearpage
\section{Fundamental Forces and the Atomic Nucleus}
\subsection{The Four Fundamental Forces and the Quest for Unification}
Modern physics recognizes \textbf{four fundamental forces} that orchestrate all interactions in nature: \textbf{gravity}, \textbf{electromagnetism}, the \textbf{strong nuclear force}, and the \textbf{weak nuclear force}. The prevailing cosmological picture, the \textbf{Big Bang}, suggests that these forces were unified in the earliest moments of the universe and differentiated as the cosmos cooled and expanded. A central objective of theoretical physics is to recover this unity through a consistent framework.

Historically, the unification of electricity and magnetism into \textbf{electromagnetism} --- captured classically by Maxwell's equations and quantum mechanically by \textbf{quantum electrodynamics (QED)} --- was the first major success. In the 20th century, electromagnetism and the weak force were unified in the \textbf{electroweak theory}, which accurately describes processes such as radioactive beta decay.

The next step toward unification is to join the \textbf{strong interaction} with the electroweak sector within a \textbf{Grand Unification Theory (GUT)}. Among candidate ideas is \textbf{supersymmetry (SUSY)}, which posits a deep relation between matter particles (fermions) and force carriers (bosons). Despite intense experimental efforts, a validated GUT remains out of reach, in part because current technology cannot probe the extremely high energies where such unification would manifest.

The most ambitious goal is a \textbf{Theory of Everything} that includes \textbf{gravity} alongside the other three forces. Proposals such as \textbf{string theory} attempt this, but they are highly theoretical and presently lack experimental confirmation.

\subsubsection{Relative Strength and Dominance Across Scales}
It is misleading to rank fundamental forces by absolute strength without specifying context. Their influence depends critically on length scale:

\begin{itemize}
\item At macroscopic and celestial scales, \textbf{gravity} dominates the dynamics of planets, stars, and galaxies, despite its small coupling at the particle level.
\item At atomic scales, \textbf{electromagnetism} governs electron-nucleus interactions, atomic structure, and all of chemistry; gravity is negligible here.
\item At nuclear scales (\ensuremath{\sim}10\ensuremath{{}^{-15}} m), the \textbf{strong} and \textbf{weak} forces play important roles. The strong force binds protons and neutrons in stable nuclei, overcoming proton-proton Coulomb repulsion. The weak force orchestrates processes such as beta decay in unstable nuclei.
\end{itemize}
The key is that each force has a characteristic range and domain of relevance: what appears ``weak'' or ``strong'' is fundamentally scale-dependent.

\subsection{The Standard Model of Particle Physics}
The \textbf{Standard Model (SM)} is the central theoretical framework that classifies all known \textbf{elementary particles} and describes their interactions via three of the four forces (gravity excluded). It is to physics what the periodic table is to chemistry: a structured catalog of constituents and rules of interaction.

\subsubsection{Matter Particles: Fermions}
All matter is built from \textbf{fermions} (spin 1/2 particles), which fall into two families:

\begin{itemize}
\item \textbf{Quarks}: Six ``flavors'' arranged in three generations --- up, down; charm, strange; top, bottom. Quarks combine to form \textbf{hadrons}:
\begin{itemize}
\item \textbf{Baryons}: three-quark composites (e.g., protons, neutrons) with spin 1/2 or 3/2.
\item \textbf{Mesons}: quark--antiquark composites (e.g., pions, kaons) with spin 0 or 1.
\end{itemize}
\end{itemize}
  A proton, for example, consists of two up quarks and one down quark.

\begin{itemize}
\item \textbf{Leptons}: The electron, muon, and tau, each with its associated neutrino. The \textbf{electron} is an elementary lepton; leptons do not experience the strong force.
\end{itemize}
For every charged matter particle, there is a corresponding \textbf{antiparticle} with the same mass and opposite electric charge.

\subsubsection{Force Carriers: Gauge Bosons}
Interactions are mediated by \textbf{gauge bosons}, the quanta exchanged between fermions:

\begin{itemize}
\item \textbf{Photon}: carrier of \textbf{electromagnetism}; massless, the quantum of light.
\item \textbf{Gluons}: carriers of the \textbf{strong force}; they ``glue'' quarks within hadrons and confine them inside nucleons.
\item \textbf{W\ensuremath{\pm} and Z\ensuremath{{}^{0}} bosons}: carriers of the \textbf{weak force}; massive and central to processes such as beta decay.
\end{itemize}
The SM does not include gravity. The hypothetical quantum of gravity, the \textbf{graviton}, has not been experimentally detected. Its absence from the SM reflects both conceptual challenges in quantizing gravity and the practical difficulty of detecting a particle associated with an extremely weak interaction at particle scales.

\subsubsection{The Higgs Boson and the Origin of Mass}
Completing the SM is the \textbf{Higgs boson}, an excitation of the \textbf{Higgs field} that permeates the universe. Elementary particles acquire \textbf{mass} through their interaction with this field. The Higgs is a \textbf{scalar} (spin-0) particle.

The Higgs boson was discovered at CERN's Large Hadron Collider by the ATLAS and CMS collaborations, supporting the Standard Model's mass-generation mechanism.

\subsection{From Atomic Models to the Quantum Atom}
Our picture of the atom evolved through successive models:

\begin{itemize}
\item \textbf{Dalton}: the atom as an indivisible solid sphere.
\item \textbf{Thomson (plum pudding)}: electrons embedded in a diffuse positive medium.
\item \textbf{Rutherford (nuclear model)}: a dense, positively charged \textbf{nucleus} with electrons orbiting it, based on the gold foil experiment.
\item \textbf{Bohr}: quantized electron orbits (energy levels), successful for single-electron systems such as hydrogen.
\item \textbf{Schrödinger (quantum mechanical model)}: electrons described by wavefunctions occupying \textbf{orbitals}, i.e., regions of probability; the state of each electron is specified by quantum numbers n, l, m\ensuremath{\ell}, m\ensuremath{{}_{s}}. This model is necessary for multi-electron atoms where Bohr's picture fails.
\end{itemize}
Bohr's model remains pedagogically useful for introducing \textbf{quantized energy levels}, but accurate atomic structure requires the full quantum treatment.

\subsection{The Atomic Nucleus: Composition, Scale, and Terminology}
The \textbf{nucleus} is the central, massive core of the atom, containing \textbf{protons} and \textbf{neutrons}, collectively called \textbf{nucleons}. It accounts for approximately 99.9\% of the atom's mass while occupying a minute fraction of its volume.

\begin{itemize}
\item \textbf{Atomic number (Z)}: the number of protons; it uniquely identifies the element and does not change for a given element.
\item \textbf{Neutron number (N)}: the number of neutrons.
\item \textbf{Mass number (A)}: the total number of nucleons, given by \textbf{A = Z + N}.
\item \textbf{Nuclear radius}: approximately obeys R \ensuremath{\approx} R\ensuremath{{}_{0}} A\ensuremath{{}^{1/3}}, with R\ensuremath{{}_{0}} \ensuremath{\approx} 1.2--1.3 fm. The nuclear scale is on the order of 10\ensuremath{{}^{-15}} m.
\end{itemize}
Standard nuclear notation writes a nuclide as \({}^{A}_{Z}X\), where X is the element symbol, A is the mass number (upper left), and Z is the atomic number (lower left). The neutron number N is often given separately (A = Z + N). In reaction equations, additional annotations (e.g., charge state q) may appear; if present, the net ionic charge is placed in the upper right.

\subsubsection{Nuclear Terminology: Isotopes, Isotones, Isobars, and Isomers}
\begin{itemize}
\item \textbf{Isotopes}: same Z, different N (thus different A). Example: \ensuremath{{}^{1}}H, \ensuremath{{}^{2}}H (deuterium), \ensuremath{{}^{3}}H (tritium).
\item \textbf{Isotones}: same N, different Z and A. Example: \ensuremath{{}^{2}}H (Z=1, N=1) and \ensuremath{{}^{3}}He (Z=2, N=1).
\item \textbf{Isobars}: same A, different Z and N. Example: \ensuremath{{}^{3}}H (Z=1, N=2) and \ensuremath{{}^{3}}He (Z=2, N=1).
\item \textbf{Isomers}: same Z, N, and A, but different nuclear energy states. A long-lived excited state is \textbf{metastable} and is denoted by ``m''. Example: \ensuremath{{}^{99m}}Tc and \ensuremath{{}^{99}}Tc (different half-lives).
\end{itemize}
A crucial principle: \textbf{All isotopes of a given element have essentially identical chemical properties} (determined by Z and the resulting electron configuration) but can have profoundly different \textbf{nuclear properties} (stability, decay modes, radiation emission), which are governed by N and nuclear structure.


\subsection{Atomic Mass Standards and the Atomic Mass Unit}
The standard atomic mass reference is \textbf{carbon-12}:

\begin{itemize}
\item \textbf{Carbon-12} is abundant and offers a stable, well-defined reference.
\item By definition, the mass of a neutral ground-state \ensuremath{{}^{12}}C atom is exactly 12 atomic mass units (amu).
\end{itemize}
One \textbf{atomic mass unit (amu or u)} is defined as one-twelfth the mass of a \ensuremath{{}^{12}}C atom:

\begin{itemize}
\item 1 u \ensuremath{\approx} 1.6605 \ensuremath{\times} 10\ensuremath{{}^{-27}} kg (more precisely, 1.66053906660 \ensuremath{\times} 10\ensuremath{{}^{-27}} kg).
\item Using rest-mass energy conversion, 1 u corresponds to approximately \textbf{931.5 MeV/c\ensuremath{{}^{2}}}.
\end{itemize}
These equivalences permit routine conversions between mass and energy via \textbf{E = mc\ensuremath{{}^{2}}}. This bridge is foundational in nuclear physics, where masses are often tabulated in energy units.

\subsubsection{Particle Masses and Rest Energies}
Key rest masses and corresponding rest energies:

\begin{itemize}
\item \textbf{Electron}: m \ensuremath{\approx} 9.11 \ensuremath{\times} 10\ensuremath{{}^{-31}} kg; E\ensuremath{{}_{0}} \ensuremath{\approx} \textbf{0.511 MeV}.
\item \textbf{Proton}: m \ensuremath{\approx} 1.67 \ensuremath{\times} 10\ensuremath{{}^{-27}} kg; E\ensuremath{{}_{0}} \ensuremath{\approx} \textbf{938 MeV}.
\item \textbf{Neutron}: m \ensuremath{\approx} 1.67 \ensuremath{\times} 10\ensuremath{{}^{-27}} kg; E\ensuremath{{}_{0}} \ensuremath{\approx} \textbf{939 MeV}.
\end{itemize}
At the nuclear scale, nucleon masses dominate energetics; at the atomic scale, electron dynamics and electromagnetism dominate structure and chemistry.


\subsection{Charge, Ionization, and Radiation}
The \textbf{atomic number (Z)} sets the element's identity. In a \textbf{neutral atom}, the number of electrons equals Z. Gain or loss of electrons produces \textbf{ions}:

\begin{itemize}
\item Gaining electrons \ensuremath{\rightarrow} negative ion (anion).
\item Losing electrons \ensuremath{\rightarrow} positive ion (cation).
\end{itemize}
Everyday electrostatics (e.g., a plastic comb attracting small paper pieces after combing hair) exemplifies charge transfer and electrostatic attraction. In radiation physics, the distinction between excitation and ionization is central:

\begin{itemize}
\item \textbf{Non-ionizing radiation} can \textbf{excite} electrons to higher energy states without removing them from atoms.
\item \textbf{Ionizing radiation} has sufficient energy to liberate electrons, producing an \textbf{ion pair}: a positive ion (the residual atom) and a free electron. The positive ion typically captures an electron from the environment to re-establish neutrality.
\end{itemize}
These processes are foundational for understanding radiation detection, biological effects, and radiation protection.

\subsection{The Chart of the Nuclides and Nuclear Stability}
Nuclear stability is systematically represented in the \textbf{Chart of the Nuclides}, which plots \textbf{neutron number (N)} versus \textbf{proton number (Z)}. Always verify axis conventions, as some resources may invert them. Stable nuclides form a narrow band often called the \textbf{line (or peninsula) of stability}. Radionuclides occupy regions away from this line, and their positions correlate with characteristic \textbf{half-lives}:

\begin{itemize}
\item In the chart discussed in class, stable isotopes are shown as black squares.
\item Long-lived unstable isotopes: red (half-lives from \textasciitilde{}350 years to \textasciitilde{}46 billion years); practically stable on human timescales.
\item Moderately unstable: yellow (half-lives 12 days to 350 years).
\item Short-lived: green (1 second to 12 days).
\item Extremely unstable: blue (less than 1 second).
\end{itemize}
An unstable nuclide tends to undergo one or more decay processes to move toward stability on the chart, changing Z and/or N accordingly.

\subsection{Nuclear Forces at the Boundary: Protons vs Neutrons}
Within the nucleus, the \textbf{strong nuclear force} is dominant and binds nucleons tightly despite proton-proton \textbf{Coulombic} repulsion. Crucially, the dominance of the strong force is confined to distances on the order of nuclear dimensions.

\begin{itemize}
\item If a \textbf{proton} is outside the nucleus, the \textbf{electromagnetic force} immediately dominates, and the positively charged proton experiences strong repulsion from the positively charged nucleus. Overcoming this \textbf{Coulomb barrier} requires substantial kinetic energy. At nuclear distances, the strong force can bind the proton.
\item A \textbf{neutron}, being electrically neutral, does not face a Coulomb barrier and can approach and enter the nucleus more readily. Nonetheless, free neutrons can be hazardous due to their mass and kinetic energy; their interactions can deposit significant energy in matter, posing serious biological risks. This underlies stringent \textbf{neutron shielding} requirements in nuclear facilities and certain medical applications.
\end{itemize}
\subsection{Mass Defect and Nuclear Binding Energy}
A central empirical fact is that the mass of a bound nucleus is less than the sum of the masses of its constituent free protons and neutrons. This \textbf{mass defect (\ensuremath{\Delta}m)} corresponds to the \textbf{binding energy (BE)} that holds the nucleus together, via \textbf{E = mc\ensuremath{{}^{2}}}. Binding energy is released when nucleons bind (exothermic), and the same energy must be supplied to disassemble the nucleus (endothermic).

Because of the MeV scale of nuclear energies, even small mass differences translate into large binding energies compared with chemical bond energies (typically eV).

To calculate BE from measured masses, one uses the conversion:

\begin{itemize}
\item 1 u \ensuremath{\approx} 931.5 MeV/c\ensuremath{{}^{2}},
\end{itemize}
so that

\begin{itemize}
\item \textbf{BE = \ensuremath{\Delta}m \ensuremath{\times} 931.5 MeV/amu}, when \ensuremath{\Delta}m is expressed in atomic mass units.
\end{itemize}
This quantitative link allows precise accounting of energy release in nuclear reactions, fission, fusion, and decay.

\subsection{Isotopes in Practice: Examples and Applications}
The lecture highlights several isotope families and uses:

\begin{itemize}
\item \textbf{Hydrogen isotopes}: \ensuremath{{}^{1}}H, \ensuremath{{}^{2}}H, \ensuremath{{}^{3}}H illustrate basic isotopic concepts and enter into fusion and tracer studies.
\item \textbf{Carbon isotopes}: \ensuremath{{}^{12}}C defines the atomic mass unit; \textbf{\ensuremath{{}^{14}}C} is radioactive and central to \textbf{radiocarbon dating}, where its known half-life and the \ensuremath{{}^{14}}C/\ensuremath{{}^{12}}C ratio enable age determinations of organic materials.
\item \textbf{Oxygen isotopes}: \ensuremath{{}^{16}}O, \ensuremath{{}^{17}}O, \ensuremath{{}^{18}}O are common and widely used in geochemistry and climatology.
\item \textbf{Cesium isotopes}: \textbf{\ensuremath{{}^{137}}Cs} is important in medicine (e.g., therapy devices and calibration sources), with a half-life of about 30.17 years.
\end{itemize}
In radionuclide production, multiple isotopes can be generated; subsequent processing must \textbf{obtain the desired radionuclide}, because chemical methods alone cannot differentiate isotopes of the same element (their chemical properties are essentially identical).


\subsection{Binding-Energy Exercises}
Calculate the nuclear binding energies of \textbf{helium-4}, \textbf{lithium-7}, \textbf{oxygen-16}, and \textbf{uranium-235} using the mass defect and the mass-energy conversion.


\clearpage
\section{Radioactive Decay and Half-Life}
\subsection{Exponential Decay and Core Equations}
\begin{itemize}
\item Decay law and differential form
\begin{itemize}
\item The rate of change of the number of radioactive atoms N with respect to time t is proportional to N: dN = -\ensuremath{\lambda} N dt, where \ensuremath{\lambda} is the decay constant (constant of proportionality).
\item Lowercase n refers to the number of atoms remaining after some time t; uppercase \(N_0\) (or N at t=0) denotes the initial number of atoms at the baseline/reference time.
\end{itemize}
\item Integration to obtain N(t)
\begin{itemize}
\item Rearrangement: (dN)/N = -\ensuremath{\lambda} dt; integrating explicitly from \(N_0\) to N and from \(t_0\) (often set to 0) to t gives ln(N) - ln(\(N_0\)) = -\ensuremath{\lambda} (t - \(t_0\)).
\item Using logarithm properties ln(N) - ln(\(N_0\)) = ln(N/\(N_0\)); exponentiating yields N/\(N_0\) = \(e^{-\lambda(t-t_0)}\).
\item Final form (with \(t_0\) = 0 for the baseline): N(t) = \(N_0\) \(e^{-\lambda t}\). If \(t_0\) is retained, the solution is N(t) = \(N_0\) \(e^{-\lambda(t-t_0)}\).
\end{itemize}
\item Activity definition and equivalence in form
\begin{itemize}
\item Activity A is the number of disintegrations per unit time; by definition A = -dN/dt = \ensuremath{\lambda} N.
\item Because A \ensuremath{\propto} N, the same exponential form applies: A(t) = \(A_0\) \(e^{-\lambda t}\), where \(A_0\) is activity at baseline time.
\end{itemize}
\end{itemize}
\subsection{Decay Constant (\ensuremath{\lambda}) and Half-Life (\(T_{1/2}\))}
\begin{itemize}
\item Unique decay constants per isotope
\begin{itemize}
\item Each radionuclide has its own decay constant, reflecting stability: more unstable isotopes have larger \ensuremath{\lambda}.
\end{itemize}
\item Deriving the \ensuremath{\lambda}--half-life relationship
\begin{itemize}
\item Half-life \(T_{1/2}\) is the time when the activity (or number of atoms) falls to half its initial value: A/\(A_0\) = 1/2 (equivalently N/\(N_0\) = 1/2).
\item Starting from A/\(A_0\) = \(e^{-\lambda T_{1/2}}\) = 1/2, take natural log: ln(1/2) = -\ensuremath{\lambda} \(T_{1/2}\) \ensuremath{\rightarrow} \ensuremath{\lambda} = ln(2) / \(T_{1/2}\).
\end{itemize}
\item Using half-life in activity equations
\begin{itemize}
\item Substituting \ensuremath{\lambda}: A(t) = \(A_0\) \(e^{-(\ln 2)t/T_{1/2}}\). This enables activity calculations when \(T_{1/2}\) is known but \ensuremath{\lambda} is not.
\end{itemize}
\end{itemize}
\subsection{Graphical Intuition and Successive Half-Lives}
\begin{itemize}
\item Ratio behavior over time
\begin{itemize}
\item At t = 0: A/\(A_0\) = 1. As time elapses, A/\(A_0\) decreases monotonically below 1 due to exponential decay.
\end{itemize}
\item First, second, third, and fourth half-lives
\begin{itemize}
\item At t = \(T_{1/2}\): A/\(A_0\) = 1/2 = 0.5.
\item At t = 2 \(T_{1/2}\): A/\(A_0\) = 1/4 = 0.25.
\item At t = 3 \(T_{1/2}\): A/\(A_0\) = 1/8 = 0.125.
\item At t = 4 \(T_{1/2}\): A/\(A_0\) = 1/16 = 0.0625.
\end{itemize}
\end{itemize}
\subsection{Asymptotic Nature of Exponential Decay}
\begin{itemize}
\item Approaching zero without reaching it
\begin{itemize}
\item The exponential decay model is asymptotic to the time axis: the predicted activity approaches zero as t increases but does not intersect the axis.
\end{itemize}
\end{itemize}
\subsection{Practical Considerations in Clinical and Regulatory Settings}
\begin{itemize}
\item Inventory and familiarity with radionuclides
\begin{itemize}
\item Because memorizing numerous \ensuremath{\lambda} and \(T_{1/2}\) values is difficult, hospitals should maintain an inventory/table of radionuclides including half-lives, decay constants, decay modes, branching percentages, and peak energies.
\item Familiarity with radionuclide specifications improves with regular use.
\end{itemize}
\item Correct classification and communication
\begin{itemize}
\item Anecdote: A radiopharmaceutical believed (from regulatory information) to be therapeutic underscored the need to consult clinicians on actual diagnostic vs therapeutic use to avoid mistakes.
\end{itemize}
\item Radiation protection implications
\begin{itemize}
\item Higher source activity corresponds to higher radiation exposure to staff. As activity decays over time, exposure decreases and may approach background radiation levels.
\item Background radiation is nonzero; a decayed check source may become close to background levels but retains some activity.
\end{itemize}
\end{itemize}
\subsection{Foundations of Radioactivity}
\begin{itemize}
\item Definition
\begin{itemize}
\item Radioactivity: spontaneous disintegration of an unstable atomic nucleus (parent) into a more stable daughter nucleus, often with higher binding energy per nucleon, accompanied by emission of energetic particles and/or electromagnetic radiation.
\end{itemize}
\item Stochastic nature and statistical laws
\begin{itemize}
\item Radioactive decay is a probabilistic (stochastic) process: individual nucleus decay times are unpredictable, but large populations follow statistical laws and branchings (percentages to various decay modes).
\end{itemize}
\item Constancy and independence of decay probability
\begin{itemize}
\item The probability of decay per unit time (\ensuremath{\lambda}) is constant for a given isotope and independent of the nuclide's age. In the basic decay model used here, ordinary changes in physical or chemical state do not alter \ensuremath{\lambda}.
\end{itemize}
\end{itemize}

\clearpage
\section{Activity and Radioactive Decay Processes}
\subsection{Activity and Units}
Activity is the number of disintegrations per unit time:

\[A(t)=-\frac{dN}{dt}=\lambda N(t).\]

The SI unit is the \textbf{becquerel (Bq)}: one disintegration per second. The legacy unit is the \textbf{curie (Ci)}, historically associated with radium-226:

\[1\,\mathrm{Bq}=1\,\mathrm{s}^{-1},\qquad 1\,\mathrm{Ci}=3.7\times10^{10}\,\mathrm{Bq}.\]

\subsection{Mean Life}
Mean life describes the average lifetime of a radioactive nucleus. It follows from the exponential probability distribution for decay:

\[p(t)=\lambda e^{-\lambda t},\qquad \tau=\frac{1}{\lambda}=\frac{T_{1/2}}{\ln 2}.\]

The discussion also considers radionuclide mixtures, using iodine-131 products containing other iodine isotopes as an example. Different radionuclides in a mixture have different decay constants and half-lives.

\subsection{Specific Activity}
\textbf{Specific activity} is activity per unit mass, expressed in Bq/g or Ci/g. \textbf{Activity concentration} can also be expressed per unit volume.

These quantities are relevant when preparing radiopharmaceuticals. The class uses fluorine-18 FDG to discuss activity relative to patient body mass, with \textbf{4 MBq/kg} as the example value. The discussion considers how patient body dimensions affect the activity needed. This example expresses administered activity per kilogram of patient body mass.

\subsection{Decay Modes and Conservation}
Radioactive decay modes reviewed include \textbf{alpha decay}, \textbf{beta-minus}, \textbf{beta-plus (positron)}, \textbf{electron capture}, \textbf{gamma emission}, \textbf{internal conversion}, \textbf{spontaneous fission}, and \textbf{proton/neutron emission}.

Across nuclear equations, check \textbf{electric charge}, \textbf{nucleon number}, \textbf{momentum}, and \textbf{energy}, including rest-mass and kinetic energy.

\subsection{Alpha Decay}
Alpha decay emits a helium-4 nucleus, decreasing the mass number by four and the atomic number by two:

\[{}^{A}_{Z}X\longrightarrow{}^{A-4}_{Z-2}Y+{}^{4}_{2}\mathrm{He}.\]

Example:

\[{}^{238}_{92}\mathrm{U}\longrightarrow{}^{234}_{90}\mathrm{Th}+{}^{4}_{2}\mathrm{He}.\]

\subsection{Beta-Minus Decay}
A neutron converts to a proton, emitting an electron and an antineutrino. The mass number is unchanged and the atomic number increases by one:

\[n\longrightarrow p+e^-+\bar\nu_e.\]

Example:

\[{}^{14}_{6}\mathrm{C}\longrightarrow{}^{14}_{7}\mathrm{N}+e^-+\bar\nu_e.\]

\subsection{Beta-Plus Decay and Positron Annihilation}
A proton converts to a neutron, emitting a positron and a neutrino. The mass number is unchanged and the atomic number decreases by one:

\[p\longrightarrow n+e^++\nu_e\quad\text{(within the nucleus).}\]

Example:

\[{}^{10}_{6}\mathrm{C}\longrightarrow{}^{10}_{5}\mathrm{B}+e^++\nu_e.\]

Positron annihilation with an electron yields two \textbf{511 keV photons}, emitted approximately \textbf{180 degrees apart}.

\subsection{Electron Capture}
An inner-shell electron is captured by the nucleus; K-shell capture is often the most probable. The process converts a proton to a neutron:

\[p+e^-\longrightarrow n+\nu_e.\]

The atomic number decreases by one while the mass number remains unchanged. Example:

\[{}^{11}_{6}\mathrm{C}+e^-\longrightarrow{}^{11}_{5}\mathrm{B}+\nu_e.\]

The inner-shell vacancy triggers \textbf{characteristic X-ray} emission as electrons fill the vacancy. These energies are characteristic of the element. The \textbf{Auger effect} is named as a related atomic relaxation process for subsequent discussion.

\subsection{Gamma Emission}
Gamma emission is nuclear de-excitation. It changes the nuclear energy state without changing the atomic number or mass number. A relatively long-lived excited state is \textbf{metastable}, denoted by \textbf{m}.

\[{}^{99\mathrm{m}}_{43}\mathrm{Tc}\longrightarrow{}^{99}_{43}\mathrm{Tc}+\gamma.\]

\subsection{Internal Conversion}
Internal conversion is an alternative nuclear de-excitation pathway involving an orbital electron. The nucleus transfers its excitation energy to that electron, which is emitted.

\subsection{Fission}
A heavy nucleus splits into lighter fragments and releases neutrons. \textbf{Spontaneous fission} occurs without an incident neutron; \textbf{neutron-induced fission} begins after neutron absorption.

The uranium-235 example illustrates neutron absorption to form excited uranium-236, followed by fission. One balanced example is:

\[{}^{235}_{92}\mathrm{U}+{}^{1}_{0}n\longrightarrow{}^{236}_{92}\mathrm{U}^{*}\longrightarrow{}^{141}_{56}\mathrm{Ba}+{}^{92}_{36}\mathrm{Kr}+3\,{}^{1}_{0}n.\]

Check the nucleon numbers: \textbf{235 + 1 = 141 + 92 + 3}. Check the proton numbers: \textbf{92 = 56 + 36}.

Released neutrons can sustain a \textbf{chain reaction}. Reactor moderators slow neutrons and affect the likelihood of subsequent fission.

\subsection{Decay and the Chart of Nuclides}
The probabilistic nature of decay governs branching and event likelihoods. Decay pathways can be related to movement on the chart of nuclides toward more stable configurations.

Nuclear stability reflects the balance between nuclear attraction and proton-proton electrical repulsion. \textbf{Even-even nuclei} are especially common among stable nuclides. Changes in proton and neutron number identify the direction associated with each decay mode.

\end{document}
